\subsection{Modulus of an automorphism}

Let G be a locally compact group, \(\varphi\) an automorphism of G, and \(\mu\) a left Haar measure on G. It is clear that \(\varphi^{-1}(\mu)\) is also a left Haar measure on G. Therefore there exists (No.~2, Th. 1) one and only one number a \textgreater{} 0 such that \(\varphi^{-1}(\mu) = a\mu\) . By No.~2, Th. 1, this number is independent of the choice of \(\mu\) . Note that if one had started with a right Haar measure, for example \(\Delta_{G}^{-1} \cdot \mu\) (No.~3), one would have arrived at the same scalar a: for, since \(\varphi^{-1}\) leaves \(\Delta_{G}\) invariant (No.~3, Remark), one has \(\varphi^{-1}(\Delta_{G}^{-1} \cdot \mu) = \Delta_{G}^{-1} \cdot \varphi^{-1}(\mu) = a\Delta_{G}^{-1} \cdot \mu\) .

DEFINITION 4. --- The number a \textgreater{} 0 such that \(\varphi^{-1}(\mu) = a\mu\) is called the modulus of the automorphism \(\varphi\) and is denoted mod G \(\varphi\) or simply mod \(\varphi\) .

If \(f\) is a \(\mu\) -integrable function on \(G\) , then

\[
\int f \big (\varphi^ {- 1} (x) \big) d \mu (x) = (\mathrm{mod} \varphi) \int f (x) d \mu (x).\tag{31}
\]

If A is a \(\mu\) -integrable subset of G, then

\[
\mu (\varphi (\mathrm{A})) = (\mathrm{mod} \varphi) \mu (\mathrm{A}).\tag{32}
\]

In particular, for \(s \in \mathbf{G}\) , let \(i_s\) be the inner automorphism \(x \mapsto s^{-1}xs\) . Then \(i_s^{-1} = \delta(s)\gamma(s)\) , therefore

\[
i _ {s} ^ {- 1} (\mu) = \delta (s) \mu = \Delta_ {\mathrm{G}} (s) \mu ,
\]

consequently

\[
\operatorname{mod} i _ {s} = \Delta_ {\mathrm{G}} (s).\tag{33}
\]

If G is either discrete or compact, then its normalized Haar measure is transformed into itself by every automorphism \(\varphi\) of G, as one sees immediately by transport of structure. Thus an automorphism of a discrete or compact group has modulus 1.

PROPOSITION 4. --- Let G be a locally compact group, \(\Gamma\) a topological group, and \(\gamma \mapsto u_{\gamma}\) a homomorphism of \(\Gamma\) into the group \(\mathcal{G}\) of automorphism of G, such that \((\gamma, x) \mapsto u_{\gamma}(x)\) is a continuous mapping of \(\Gamma \times G\) into G. Then, the mapping \(\gamma \mapsto \operatorname{mod}(u_{\gamma})\) is a continuous representation of \(\Gamma\) in \(\mathbf{R}_{+}^{*}\) .

This mapping is obviously a representation (algebraic) of \(\Gamma\) in \(\mathbf{R}_+^*\) ; it will suffice to prove its continuity. Let \(f \in \mathcal{K}(\mathbf{G})\) and let \(\mathbf{S}\) be its support.

Let \(\gamma_0 \in \Gamma\) and let \(U\) be a relatively compact neighborhood of \(u_{\gamma_0}^{-1}(S)\) . The mapping \(\gamma \mapsto u_\gamma\) is a continuous mapping of \(\Gamma\) into \(\mathcal{G}\) equipped with the topology of compact convergence (GT, X, §3, No.~4, Th. 3); therefore \(u_{\gamma}^{-1}(S) \subset U\) for \(\gamma\) sufficiently near \(\gamma_0\) . Lemma 1 of No.~1 then proves that \(\int f(u_\gamma(x)) d\mu(x)\) (where \(\mu\) denotes a left Haar measure of \(G\) ) depends continuously on \(\gamma\) ; whence the proposition.

